% ==================================================================
% id:       lim-039
% topic:    גבולות ורציפות של פונקציות
% subtopic: משפט הסנדוויץ׳
% solution: full
% ==================================================================
\textbf{הוכיחו או הפריכו כל אחד מהסעיפים הבאים:}
\begin{enumerate}
\item אם $|f(x)-3|\le|x-1|$ לכל $x\ne1$, אז $\ds\lim_{x\to1}f(x)=3$.
\ifsol
\begin{solbox}
הטענה נכונה. לכל $x\ne1$ מתקיים
\[3-|x-1|\le f(x)\le3+|x-1|\]
ושני החסמים שואפים ל-$3$. לפי הסנדוויץ' $\ds\lim_{x\to1}f(x)=3$.
\end{solbox}
\fi
\item אם $h(x)\le f(x)\le g(x)$ בסביבה נקובה של $x_0$, וקיימים הגבולות של $h$ ושל $g$ ב-$x_0$, אז קיים הגבול של $f$ ב-$x_0$.
\ifsol
\begin{solbox}
הטענה אינה נכונה. למשל $x_0=1$, $h(x)=0$, $g(x)=1$ ו-$f(x)=\floor x$. בקטע $(0,2)$ מתקיים $0\le\floor x\le1$, אבל ל-$\floor x$ אין גבול ב-$1$: עבור $1-\frac1n$ מתקבל $0$, ועבור $1+\frac1n$ מתקבל $1$.
\end{solbox}
\fi
\item אם $\ds\lim_{x\to0}f(x)=0$, ו-$g$ מוגדרת בסביבה נקובה של $0$, אז $\ds\lim_{x\to0}f(x)g(x)=0$.
\ifsol
\begin{solbox}
הטענה אינה נכונה. למשל $f(x)=x$ ו-$g(x)=\frac1x$: אז $f(x)g(x)=1$ לכל $x\ne0$, והגבול הוא $1$. ``אפסה כפול חסומה'' דורש ש-$g$ תהיה חסומה.
\end{solbox}
\fi
\item אם $\ds\lim_{x\to x_0}|f(x)|=0$, אז $\ds\lim_{x\to x_0}f(x)=0$.
\ifsol
\begin{solbox}
הטענה נכונה. לכל $x$ מתקיים $-|f(x)|\le f(x)\le|f(x)|$, ולפי הנתון $\ds\lim_{x\to x_0}\lr{-|f(x)|}=\lim_{x\to x_0}|f(x)|=0$. לפי משפט הסנדוויץ' $\ds\lim_{x\to x_0}f(x)=0$.
\end{solbox}
\fi
\item אם $\ds\lim_{x\to x_0}|f(x)|=1$, אז $\ds\lim_{x\to x_0}f(x)=1$ או $\ds\lim_{x\to x_0}f(x)=-1$.
\ifsol
\begin{solbox}
הטענה אינה נכונה. למשל $x_0=0$ ו-$f(x)=1$ עבור $x\ge0$, ו-$f(x)=-1$ עבור $x<0$. אז $|f(x)|=1$ לכל $x$, אבל ל-$f$ אין גבול ב-$0$: עבור $\frac1n$ מתקבל $1$, ועבור $-\frac1n$ מתקבל $(-1)$.
\end{solbox}
\fi
\end{enumerate}
